\section{Certifying an instance edition and a route result}\label{app:certification-order}

Consider a family that distributes travel times for a set of customer locations. A team submits a shorter tour a month after the instance is published. The new tour changes the incumbent, while the locations, travel times and rules under which the two tours are compared stay fixed. An instance edition therefore has its own identity, which the family can publish before it receives a solution.

\paragraph{Publish the inputs first.}
The publisher assigns an edition label such as \id{r1} to an instance variant within a family. The descriptor and required sidecars are available as files with content hashes. For regenerated operands, the publisher also preserves the source distribution containing the materialization function, its parameters and the relevant environment information. These artifacts must remain retrievable by their recorded digests. The complete matrix or ATF stream is then computed and hashed as specified in Appendix~\ref{app:materialization-format}. If the family supplies both a stored representation and a regeneration procedure for the same arc data, a verifier can compare the two stream digests directly.

\paragraph{Publish the manifest second.}
For a descriptor named \path{C101.vrp.json}, its edition manifest is \path{C101.manifest.json} in the same directory. The manifest records the edition label in \id{benchmark_revision}, the objective and numerical profile, the input references, and the expected materialization digest. Once these fields are fixed, the publisher preserves the exact manifest bytes. Changing an input, the objective or the prescribed recipe requires a new edition of that variant. Previously issued manifests and their inputs must remain available for replay.

\paragraph{Bind the result last.}
The verifier reads the manifest, checks the input and materialization digests, and evaluates the submitted routes under the declared objective and profile. The result records the edition label and the SHA-256 of the complete manifest file in \id{metadata.standard}. The digest includes the manifest's whitespace and final newline. It appears in the result, outside the bytes being hashed. The manifest refers only to inputs and materialization artifacts, so replacing a BKS or verifying it with another conforming backend leaves the manifest intact. Historical solver claims identify their own target as described in Appendix~\ref{app:solver-claims}.

\paragraph{Verification after a repository move.}
A Git commit identifies a repository snapshot together with its history. Replaying the same files on a different parent can change the commit identifier while preserving each file byte for byte. A release archive may contain the same files without any Git metadata. In both cases the verifier can follow the result's manifest digest, check the descriptor and sidecars by their file digests, and reproduce the matrix or ATF digest. A version-control reference can help a reader locate a source. Certification depends on the preserved artifact bytes. For example, moving a copied instance and its manifest into a directory outside a checkout preserves their identity. Editing a single customer deadline in that copy invalidates the descriptor hash even if every entry in the travel matrix remains identical.
